\documentclass{article}%
\usepackage[T1]{fontenc}%
\usepackage[utf8]{inputenc}%
\usepackage{lmodern}%
\usepackage{textcomp}%
\usepackage{lastpage}%
\usepackage{geometry}%
\geometry{margin=2.2cm}%
%
\usepackage{amsmath}%
\usepackage{amssymb}%
\usepackage{booktabs}%
\usepackage{longtable}%
\usepackage{array}%
\usepackage{hyperref}%
\hypersetup{hidelinks}%
\title{IME 2026 {-} Química (claude\_sonnet\_5\_5) x Poliedro\\[4pt]\large Comparação com o gabarito Poliedro Resolve\thanks{Gerado com gpt-resolve: \url{https://github.com/lgabs/gpt-resolve}}}%
\date{10/11/2026}%
%
\begin{document}%
\normalsize%
\maketitle%
\section*{Resumo}\textbf{Nota média: 10,0/10} (16 de 16 itens com a resolução correta; cada item vale o mesmo dentro da questão).\par\medskip%
\begin{longtable}{@{}c c p{0.78\textwidth}@{}}\toprule Questão & Nota & Resumo \\ \midrule
1 & 10,0 & A resolução está correta e bate com o gabarito do Poliedro: Q = 9648 C = 0,1 F, n(M) = 0,05 mol, massa molar 87,6 g/mol, e o metal é o estrôncio (Sr). Refiz as contas e conferi a massa do Sr na tabela da prova (87,62). Há só uma observação de redação, sobre a parte de incerteza, que não afeta a nota. \\
2 & 10,0 & O item a coincide: C é um álcool terciário. No item b, B (2{-}metilbut{-}1{-}eno) e C (2{-}metilbutan{-}2{-}ol) também coincidem. Resolução e gabarito divergem só em A. A resolução escolhe o 1{-}cloro{-}2{-}metilbutano, que gera B sem mistura; o Poliedro escolhe o 2{-}cloro{-}2{-}metilbutano, que gera B como produto principal só com base volumosa (Hofmann) e, com base comum (Zaitsev), dá principalmente o 2{-}metilbut{-}2{-}eno. Como o enunciado não fixa a base, as duas respostas são quimicamente válidas, e a da resolução é a mais rigorosa; ela ainda cita a do gabarito como alternativa. O risco de nota está em A, caso a banca aceite só o haleto terciário. \\
3 & 10,0 & A resolução e o gabarito do Poliedro concordam nos dois itens, e ambos estão corretos. a) A fórmula molecular é C12H20O2, e a estrutura é CH2=CH--C(CH3)(OAc)--CH2--CH2--CH=C(CH3)2 (acetato de 3,7{-}dimetilocta{-}1,6{-}dien{-}3{-}ila); a contagem de átomos e o índice de insaturação conferem. b) A poliadição ocorre só na dupla trissubstituída C6=C7 e dá a unidade repetitiva {[}--CH(R)--C(CH3)2--{]}n, com R = --CH2CH2C(CH3)(OAc)CH=CH2 e o vinil livre. É o mesmo polímero do gabarito, escrito na ordem inversa. Conferi o enunciado no PDF original, e ele confere com a transcrição. Restam só observações de redação e de justificativa, nenhuma capaz de mudar a nota. \\
4 & 10,0 & A resolução e o gabarito do Poliedro chegam ao mesmo resultado: são necessários 120 mol de ácido acético por piscina, cerca de 7,2 kg. A resolução usa o balanço de cargas e de massa, e o gabarito usa a Lei de Ostwald sem aproximação, com $\alpha$ = 1/6. Os dois caminhos são equivalentes. Refiz as contas e os dois estão certos. As únicas ressalvas são hipóteses implícitas: a água não tem efeito tampão e o metal não interfere no equilíbrio do ácido. \\
5 & 10,0 & A resolução e o gabarito Poliedro dão a mesma resposta: m1/m2 = 1 (m1 = m2). Refiz a conta e o resultado está certo. Como a energia é proporcional à quantidade de reação, a 1ª queima consome 4 vezes mais CH4 que a 2ª. Esse fator 4 é compensado pela razão entre as massas, 2$\cdot$32 g de O2 para 16 g de CH4. Com as massas atômicas da tabela, a razão dá 1,003, o que é só arredondamento. Não há erro em nenhum dos dois, e os únicos pontos de dúvida são menores, de redação e de hipótese implícita. \\
6 & 10,0 & A resolução e o gabarito Poliedro dão a mesma resposta e os dois estão certos: o ácido é o etanoico (acético), CH3--COOH (C2H4O2, M = 60 g/mol, monoácido). Conferi o enunciado no PDF e as contas em Python: n(NaOH) = 3,0$\times$10\^{}{-}3 mol, M = 60 g/mol, k = 1 e C2H4O2 como única fórmula possível. Os caminhos são um pouco diferentes: a resolução tira o radical pela massa (60 - 45 = 15, metil) e o gabarito resolve 12x + y = 28. A resolução é mais completa, porque exclui o diácido e confere o resultado com as massas da tabela da prova. Ficam só duas observações de baixa importância: a hipótese de dissociação completa não vale de verdade para um ácido fraco, e o vapor de ácido acético forma dímeros. Nenhuma das duas muda a nota. \\
7 & 10,0 & A resolução está certa e bate com o gabarito do Poliedro: 42 placas. As contas conferem: $\Delta$H\_fus(Cu) = 12,9648 kJ/mol, $\Delta$H global = +49,93 kJ por 2 mol de Cu, n\_Cu = 3357,6 mol, Q = 83,82 MJ e N = 41,91, logo 42. A diferença entre as duas é só de caminho para chegar a $\Delta$H\_fus: a resolução identifica os coeficientes de $\Delta$G\_fus, e o gabarito impõe $\Delta$G = 0 em T\_f. O resultado é o mesmo. Para levar a um professor: o enunciado diz kJ/mol nas expressões de G, mas os dados só fazem sentido em J/mol. Os dois adotam J/mol, e a resolução diz isso explicitamente. \\
8 & 10,0 & A resolução está correta e coincide com o gabarito Poliedro: 50 elétrons reais contra 84 esperados (30 B--B + 12 B--H em ligações 2c--2e), deficiência de 34/84 $\approx$ 40,5\%. Conferi as contas em Python e o enunciado contra o PDF original. A resolução vai além do gabarito com a verificação por Euler/3F = 2A, a comparação com o B2H6 e a regra de Wade, e discute e descarta de forma explícita as leituras alternativas (68\% e 56,7\%). Não há erro a apontar. A única ressalva é que o enunciado não define a fórmula da deficiência. \\
9 & 10,0 & Questão 9: a resolução bate com o gabarito do Poliedro nos três itens, e as duas estão certas. Refiz as contas em Python: (a) {[}CO2{]} = 1,8$\times$10\^{}{-}5 mol/L; (b) $\Delta$G$^\circ$ = +36,3 kJ/mol, K $\approx$ 4,2$\times$10\^{}{-}7, pH $\approx$ 5,56 $\approx$ 5,6; (c) $\Delta$H$^\circ$ = +7,6 kJ/mol, ionização endotérmica, então o pH diminui. A resolução se justifica melhor que o gabarito: confere as aproximações (Kw e 2ª ionização) e estima o tamanho do efeito. O gabarito só tem erros de digitação que não mudam nada (\textbackslash{}"CO3\^{}{-}\textbackslash{}" no lugar de \textbackslash{}"HCO3\^{}{-}\textbackslash{}" e uma equação truncada). A única hipótese que vale citar é o uso da constante de Henry de 300 K a 25 $^\circ$C, adotada pelos dois. \\
10 & 10,0 & A resolução da questão 10 bate com o gabarito do Poliedro nos três itens. a) As etapas são as mesmas ($\beta$$^{-}$ e depois três $\alpha$), com balanço correto. b) Família do actínio (4n+3). c) $\lambda$ $\approx$ 0,173 s$^{-}$$^{1}$ e $\tau$ $\approx$ 5,77 s; a resolução arredonda para 0,17 s$^{-}$$^{1}$ e 5,8 s. Ambos estão certos. A observação do gabarito sobre um suposto \textbackslash{}"277\textbackslash{}" no enunciado não vale para o PDF original desta pasta, que traz 227. \\
\bottomrule\end{longtable}%
\section*{Principais pontos de dúvida}%
\begin{itemize}%
\item%
\textbf{Q2} \textit{(key)} O A do gabarito (2{-}cloro{-}2{-}metilbutano, haleto terciário) dá 2{-}metilbut{-}1{-}eno apenas como produto minoritário (Hofmann), a menos que se use base volumosa. Com base comum (KOH/etanol, Zaitsev), o produto principal é o 2{-}metilbut{-}2{-}eno, que não gera formaldeído na ozonólise. O 1{-}cloro{-}2{-}metilbutano, escolhido pela resolução, gera B sem mistura. Vale perguntar ao professor se a banca aceitaria os dois haletos ou só um deles. Se o critério de correção seguir o Poliedro à risca, a resolução pode perder parte do item b, mesmo tendo a resposta quimicamente mais rigorosa.%
\end{itemize}%
\section*{Questão 1}%
\href{https://poliedroresolve.sistemapoliedro.com.br/vestibulares/ime/2026/ime-2-fase-dia-3-25-10-2025-2-fase-dia-3/dissertativa-1-quimica-2-fase-dia-3-ime-2026}{Resolução do gabarito}\par\smallskip%
\begin{longtable}{@{}p{0.06\textwidth} p{0.33\textwidth} p{0.33\textwidth} p{0.1\textwidth} p{0.08\textwidth}@{}}\toprule Item & Resolução & Gabarito & Veredito & Certo \\ \midrule
Questão 1 -- identificação do metal M & M = estrôncio (Sr, Z = 38): Q = 9648 C, n(e$^{-}$) = 0,100 mol, n(M) = 0,050 mol, massa molar 87,6 g/mol & Q = 9648 C = 0,1 F; M$^2$? + 2e$^{-}$ $\to$ M; massa molar 87,6 g/mol $\Rightarrow$ M é o estrôncio (Sr) & igual & ambos \\
\bottomrule\end{longtable}%
\begin{itemize}%
\item%
\textbf{Questão 1 -- identificação do metal M)} O enunciado transcrito confere com o original (página 2). Refiz as contas: Q = 0,4$\cdot$6,7$\cdot$3600 = 9648 C; com F = 96480 C/mol, n(e$^{-}$) = 0,1 mol, e como z = 2, n(M) = 0,05 mol. Então M = 4,38/0,05 = 87,6 g/mol. A tabela da página 1 dá Sr = 87,62, Rb = 85,468 e Y = 88,906, o que confirma o Sr. A resolução e o gabarito chegam ao mesmo resultado pelo mesmo caminho. As verificações extras da resolução também estão certas: com F = 96500 sai 87,62; com z = 1 ou 3 sairia 43,8 ou 131,4 g/mol; a incerteza propagada de t e m dá cerca de $\pm$0,66 a 0,75 g/mol.%
\end{itemize}%
\paragraph{Pontos de dúvida}%
\begin{itemize}%
\item%
\textit{(solution)} A discussão de incerteza no item 7 não é necessária e tem um argumento fraco. Ela deixa de fora a incerteza de i = 0,4 A, que tem 1 algarismo significativo, alegando que 'só um valor nominal exato permite identificar o elemento'. Se essa incerteza fosse levada a sério, o resultado não permitiria distinguir os elementos vizinhos. O correto é dizer apenas que os dados são nominais, já que Q = 0,1 F exatamente. Isso não afeta a resposta nem a nota.%
\end{itemize}%
\section*{Questão 2}%
\href{https://poliedroresolve.sistemapoliedro.com.br/vestibulares/ime/2026/ime-2-fase-dia-3-25-10-2025-2-fase-dia-3/dissertativa-2-quimica-2-fase-dia-3-ime-2026}{Resolução do gabarito}\par\smallskip%
\begin{longtable}{@{}p{0.06\textwidth} p{0.33\textwidth} p{0.33\textwidth} p{0.1\textwidth} p{0.08\textwidth}@{}}\toprule Item & Resolução & Gabarito & Veredito & Certo \\ \midrule
a & Álcool (terciário): C não é oxidado por KMnO4/H+ porque o carbono que leva o OH não tem H. & Álcool (terciário, pois não reage com KMnO4 em meio ácido). & igual & ambos \\
b & A = 1{-}cloro{-}2{-}metilbutano, ClCH2CH(CH3)CH2CH3 (cita o 2{-}cloro{-}2{-}metilbutano como alternativa, com base volumosa); B = 2{-}metilbut{-}1{-}eno, CH2=C(CH3)CH2CH3; C = 2{-}metilbutan{-}2{-}ol, (CH3)2C(OH)CH2CH3. & A = 2{-}cloro{-}2{-}metilbutano, CH3CH2C(CH3)2Cl; B = 2{-}metilbut{-}1{-}eno; C = 2{-}metilbutan{-}2{-}ol. & diferente & ambos \\
\bottomrule\end{longtable}%
\begin{itemize}%
\item%
\textbf{a)} Mesma função e mesma justificativa: álcoois primários e secundários são oxidados por KMnO4/H+, o terciário não. Conferi o enunciado no PDF original (pág. 2) e a transcrição está fiel.%
\item%
\textbf{b)} B e C coincidem. A diferença está só em A. Refiz a análise: sobre o esqueleto do 2{-}metilbutano, só dois cloretos dão 2{-}metilbut{-}1{-}eno por eliminação de HCl. (i) 1{-}cloro{-}2{-}metilbutano: o único C{-}$\beta$ é o C2, então B é o único alceno formado. (ii) 2{-}cloro{-}2{-}metilbutano: os C{-}$\beta$ são C1, C3 e a metila, e formam{-}se 2{-}metilbut{-}1{-}eno e 2{-}metilbut{-}2{-}eno. Com base pequena (Zaitsev), o principal é o 2{-}metilbut{-}2{-}eno, que na ozonólise dá propanona e etanal, sem formaldeído. Com base volumosa (Hofmann), o principal é B. Os outros dois cloretos possíveis não servem: o 1{-}cloro{-}3{-}metilbutano dá 3{-}metilbut{-}1{-}eno, e o 2{-}cloro{-}3{-}metilbutano dá 2{-}metilbut{-}2{-}eno ou 3{-}metilbut{-}1{-}eno. Como o enunciado não diz qual base é usada nem exige o produto principal, as duas respostas para A cabem. A da resolução é a mais rigorosa porque forma B sem mistura. A do gabarito também forma B, mas só como produto minoritário quando a base é pequena. Por isso marquei as duas como corretas. A resolução, aliás, já aponta a resposta do gabarito como alternativa válida.%
\end{itemize}%
\paragraph{Pontos de dúvida}%
\begin{itemize}%
\item%
\textit{(key)} O A do gabarito (2{-}cloro{-}2{-}metilbutano, haleto terciário) dá 2{-}metilbut{-}1{-}eno apenas como produto minoritário (Hofmann), a menos que se use base volumosa. Com base comum (KOH/etanol, Zaitsev), o produto principal é o 2{-}metilbut{-}2{-}eno, que não gera formaldeído na ozonólise. O 1{-}cloro{-}2{-}metilbutano, escolhido pela resolução, gera B sem mistura. Vale perguntar ao professor se a banca aceitaria os dois haletos ou só um deles. Se o critério de correção seguir o Poliedro à risca, a resolução pode perder parte do item b, mesmo tendo a resposta quimicamente mais rigorosa.%
\item%
\textit{(statement)} Ambiguidade do enunciado em B: na hidratação ácida do 3{-}metilbut{-}1{-}eno, o carbocátion secundário sofre migração 1,2 de hidreto e vira terciário, e na prática o produto principal é o 2{-}metilbutan{-}2{-}ol (álcool terciário). Assim, o 3{-}metilbut{-}1{-}eno também seria compatível com C, e A poderia ser o 1{-}cloro{-}3{-}metilbutano. O gabarito despreza o rearranjo e diz que esse alceno dá o 3{-}metilbutan{-}2{-}ol. A resolução reconhece o rearranjo, mas adota a mesma escolha do gabarito (Markovnikov direto). B final é o mesmo nos dois.%
\item%
\textit{(solution)} Na resposta final, a resolução dá uma alternativa para A ("2{-}cloro{-}2{-}metilbutano com base volumosa"). A condição está explícita e o argumento está certo, mas alguns corretores penalizam respostas com mais de uma opção. A resposta principal (1{-}cloro{-}2{-}metilbutano) está bem destacada.%
\item%
\textit{(key)} Segundo a transcrição, na imagem da resolução do gabarito o 3{-}metilbutan{-}2{-}ol aparece desenhado com uma dupla ligação junto ao carbono do OH. É erro de desenho e não afeta a resposta final.%
\end{itemize}%
\section*{Questão 3}%
\href{https://poliedroresolve.sistemapoliedro.com.br/vestibulares/ime/2026/ime-2-fase-dia-3-25-10-2025-2-fase-dia-3/dissertativa-3-quimica-2-fase-dia-3-ime-2026}{Resolução do gabarito}\par\smallskip%
\begin{longtable}{@{}p{0.06\textwidth} p{0.33\textwidth} p{0.33\textwidth} p{0.1\textwidth} p{0.08\textwidth}@{}}\toprule Item & Resolução & Gabarito & Veredito & Certo \\ \midrule
a & CH2=CH--C(CH3)(OCOCH3)--CH2--CH2--CH=C(CH3)2 (acetato de 3,7{-}dimetilocta{-}1,6{-}dien{-}3{-}ila, C12H20O2) & CH2=CH--C(CH3)(OCOCH3)--CH2--CH2--CH=C(CH3)2 = CH3COO--C(CH3)(CH=CH2)--CH2CH2CH=C(CH3)2 (acetato de 3,7{-}dimetilocta{-}1,6{-}dien{-}3{-}ila, C12H20O2) & igual & ambos \\
b & Poliadição na dupla trissubstituída C6=C7: {[}--CH(R)--C(CH3)2--{]}n, com R = --CH2CH2C(CH3)(OCOCH3)CH=CH2 (o vinil e o acetato ficam pendentes) & --{[}C(CH3)2--CH(CH2CH2C(CH3)(OCOCH3)CH=CH2){]}n--, com o vinil C1=C2 livre & equivalente & ambos \\
\bottomrule\end{longtable}%
\begin{itemize}%
\item%
\textbf{a)} A estrutura é a mesma. Conferi o raciocínio. A fórmula mínima C6H10O com 2 O dá C12H20O2, com IDH = (2$\cdot$12+2-20)/2 = 3, ou seja, duas C=C e uma C=O. A cadeia principal de 8 C tem duplas em 1 e 6, separadas por C3--C5, logo não conjugadas. As duas metilas da mesma ponta são C8 e a ramificação em C7. Os carbonos alílicos não terminais são C3 e C5; C3 fica ao lado do vinil, que é menos impedido. Por isso o C3 recebe a metila e o acetato. A contagem de H dá 2+1+3+3+2+2+1+6 = 20. A estrutura é a do acetato de linalila real e coincide com a imagem do gabarito.%
\item%
\textbf{b)} As duas respostas têm a mesma unidade repetitiva, só escrita na ordem inversa (CH(R)--C(CH3)2 contra C(CH3)2--CH(R)), e isso não muda o polímero. Nas duas, a ligação $\pi$ de C6=C7 (a dupla mais substituída) se abre, a cadeia principal fica formada por --CH(R)--C(CH3)2--, e o vinil e o acetato ficam como grupos laterais. A unidade repetitiva tem a mesma fórmula do monômero, C12H20O2. O gabarito escreve a equação n monômero $\to$ polímero e a resolução mostra só o polímero, mas o enunciado pede apenas 'o polímero possivelmente formado'.%
\end{itemize}%
\paragraph{Pontos de dúvida}%
\begin{itemize}%
\item%
\textit{(statement)} O critério (v), 'carbono alílico não terminal de menor impedimento espacial', é qualitativo. O gabarito justifica a escolha de C3 por ele estar 'mais afastado do grupo metil'; a resolução, por estar vizinho ao vinil monossubstituído, que é menos congestionado que a dupla trissubstituída perto de C5. As duas justificativas levam a C3, e a estrutura resultante é a do acetato de linalila real. Ainda assim, um corretor pode aceitar mais de uma forma de justificar.%
\item%
\textit{(solution)} No item b, a resolução não escreve a equação n monômero $\to$ polímero, como faz o gabarito; apresenta só o polímero. Isso basta para o que o enunciado pede. Ela também não discute a regiorregularidade (cabeça{-}cauda) do encadeamento, mas o gabarito também não, e para a unidade repetitiva isso é irrelevante.%
\item%
\textit{(solution)} A resolução observa que alcenos trissubstituídos polimerizam mal na prática e que a iniciação catiônica seria favorecida pelo carbocátion terciário em C7. Esses comentários estão quimicamente corretos e não são pedidos. Eles não contradizem o gabarito, que trata a polimerização apenas como hipótese ('possivelmente formado').%
\end{itemize}%
\section*{Questão 4}%
\href{https://poliedroresolve.sistemapoliedro.com.br/vestibulares/ime/2026/ime-2-fase-dia-3-25-10-2025-2-fase-dia-3/dissertativa-4-quimica-2-fase-dia-3-ime-2026}{Resolução do gabarito}\par\smallskip%
\begin{longtable}{@{}p{0.06\textwidth} p{0.33\textwidth} p{0.33\textwidth} p{0.1\textwidth} p{0.08\textwidth}@{}}\toprule Item & Resolução & Gabarito & Veredito & Certo \\ \midrule
Q4 -- massa mínima de ácido acético por piscina & $\approx$ 7,2$\times$10$^3$ g = 7,2 kg (120 mol), com pH = 4, {[}HA{]} = 5,0$\times$10$^{-}$$^{4}$ mol/L, C = 6,0$\times$10$^{-}$$^{4}$ mol/L & 7200 g = 7,2 kg (120 mol), com $\alpha$ = 1/6 e M = 6$\times$10$^{-}$$^{4}$ mol/L & igual & ambos \\
\bottomrule\end{longtable}%
\begin{itemize}%
\item%
\textbf{Q4 -- massa mínima de ácido acético por piscina)} Refiz a conta. Com {[}H?{]} = 10$^{-}$$^{4}$ mol/L (pH = 4), {[}A$^{-}${]} $\approx$ 10$^{-}$$^{4}$ mol/L e {[}HA{]} = (10$^{-}$$^{4}$)$^2$/(2$\times$10$^{-}$$^{5}$) = 5$\times$10$^{-}$$^{4}$ mol/L. Então C = 6$\times$10$^{-}$$^{4}$ mol/L, n = 6$\times$10$^{-}$$^{4}$ $\times$ 2$\times$10$^{5}$ = 120 mol e m $\approx$ 120 $\times$ 60 = 7200 g. Com M = 60,05 g/mol, m $\approx$ 7206 g; a diferença é só de arredondamento. A resolução usa o balanço de massa e de cargas, e o gabarito usa a Lei de Ostwald sem aproximação ($\alpha$ = 1/6). Os dois métodos são equivalentes e dão o mesmo valor. A transcrição do enunciado confere com o PDF.%
\end{itemize}%
\paragraph{Pontos de dúvida}%
\begin{itemize}%
\item%
\textit{(statement)} As duas soluções supõem que a água de rejeito não tem capacidade tampão e que o sal do metal pesado (por hidrólise ou pela própria precipitação) não consome nem libera H?. O enunciado não diz isso. Se a precipitação consumisse ácido, a massa mínima seria maior. É a hipótese que se espera no vestibular, mas fica implícita.%
\item%
\textit{(key)} O gabarito começa com Ka = M$\alpha$$^2$ e obtém $\alpha$ = 0,2. Esse valor só serve para mostrar que a aproximação não vale, porque $\alpha$ > 10\%. Em seguida o gabarito usa a expressão exata e chega a $\alpha$ = 1/6. O resumo do gabarito cita '$\alpha$ = 0,2', mas o valor final correto é $\alpha$ = 1/6 $\approx$ 16,7\%.%
\item%
\textit{(solution)} A resolução calcula M = 60,05 g/mol, o que daria cerca de 7206 g, mas arredonda o resultado para 7,2$\times$10$^3$ g. O gabarito usa M = 60 g/mol. A diferença é só de arredondamento.%
\end{itemize}%
\section*{Questão 5}%
\href{https://poliedroresolve.sistemapoliedro.com.br/vestibulares/ime/2026/ime-2-fase-dia-3-25-10-2025-2-fase-dia-3/dissertativa-5-quimica-2-fase-dia-3-ime-2026}{Resolução do gabarito}\par\smallskip%
\begin{longtable}{@{}p{0.06\textwidth} p{0.33\textwidth} p{0.33\textwidth} p{0.1\textwidth} p{0.08\textwidth}@{}}\toprule Item & Resolução & Gabarito & Veredito & Certo \\ \midrule
Razão m1/m2 & m1/m2 = 1 (m1 = m2); $\approx$ 1,003 com as massas atômicas da tabela & m1/m2 = 1 (m1 = m2) & igual & ambos \\
\bottomrule\end{longtable}%
\begin{itemize}%
\item%
\textbf{Razão m1/m2)} A transcrição confere com o original (p. 3). Conta refeita: o calor liberado a pressão constante é proporcional à extensão da reação CH4 + 2 O2 $\to$ CO2 + 2 H2O(l). Logo n(CH4) na 1ª queima = 4 $\times$ n(CH4) na 2ª (324/81). Na 2ª queima, n(O2) = 2$\cdot$n2(CH4). Então m1/m2 = (4$\cdot$16)/(2$\cdot$32) = 1. Com as massas da tabela (16,0426 e 31,998), dá 1,0027 $\approx$ 1. Com $\Delta$H $\approx$ -890 kJ/mol, m1 $\approx$ 5,84 g e m2 $\approx$ 5,82 g, o que confirma o resultado. O gabarito chega ao mesmo valor por outro caminho (y mol de O2 para 81 kJ, 4y mol para 324 kJ, x = 2y mol de CH4), e o raciocínio está certo.%
\end{itemize}%
\paragraph{Pontos de dúvida}%
\begin{itemize}%
\item%
\textit{(statement)} 'Determine a razão entre m1 e m2' não diz se é m1/m2 ou m2/m1. Aqui isso não importa, porque a razão vale 1 nos dois sentidos.%
\item%
\textit{(solution)} A resolução e o gabarito supõem que 'nas mesmas condições' também quer dizer água líquida na 2ª queima, ou seja, o mesmo $\Delta$H por mol de reação. Essa é a leitura natural do enunciado. Os dois usam M(CH4) = 16 e M(O2) = 32. Com a tabela da prova, a razão dá 1,003, uma diferença só de arredondamento.%
\end{itemize}%
\section*{Questão 6}%
\href{https://poliedroresolve.sistemapoliedro.com.br/vestibulares/ime/2026/ime-2-fase-dia-3-25-10-2025-2-fase-dia-3/dissertativa-6-quimica-2-fase-dia-3-ime-2026}{Resolução do gabarito}\par\smallskip%
\begin{longtable}{@{}p{0.06\textwidth} p{0.33\textwidth} p{0.33\textwidth} p{0.1\textwidth} p{0.08\textwidth}@{}}\toprule Item & Resolução & Gabarito & Veredito & Certo \\ \midrule
Nome e estrutura molecular do ácido & Ácido etanoico (ácido acético), CH3--COOH, ou seja, H3C--C(=O)--OH (C2H4O2, M = 60 g/mol, monocarboxílico; n(NaOH) = 3,0$\times$10\^{}{-}3 mol) & Ácido etanoico (ácido acético), C2H4O2 (M = 60 g/mol, monoácido): CH3--COOH, ou seja, H3C--C(=O)--OH & igual & ambos \\
\bottomrule\end{longtable}%
\begin{itemize}%
\item%
\textbf{Nome e estrutura molecular do ácido)} Conferi o enunciado na página 3 do PDF: está igual à transcrição. As contas batem: n(NaOH) = 0,1 $\times$ 0,030 = 3,0$\times$10\^{}{-}3 mol; M = 30 $\times$ 2 = 60 g/mol; k$\cdot$0,18/60 = 3$\times$10\^{}{-}3 dá k = 1, então é um monoácido. O gabarito chega à fórmula por 12x + y = 28 com O2: a única saída possível é x = 2, y = 4, porque x = 1 exigiria y = 16. A resolução chega ao mesmo resultado por M(R) = 60 - 45 = 15, que é o metil. As duas dão o mesmo nome (ácido etanoico/acético) e a mesma fórmula estrutural, H3C--C(=O)--OH. A verificação final da resolução (0,18/60,05 $\approx$ 3,0$\times$10\^{}{-}3 mol; 60,05/2,016 $\approx$ 29,8 $\approx$ 30) também confere.%
\end{itemize}%
\paragraph{Pontos de dúvida}%
\begin{itemize}%
\item%
\textit{(statement)} A hipótese de 'dissociação completa' não é realista para o ácido acético, que é um ácido fraco (Ka $\approx$ 1,8$\times$10\^{}{-}5). No enunciado ela só serve para dizer que cada --COOH é neutralizado por um OH\^{}{-} (estequiometria 1:1 até o ponto de equivalência). Além disso, o vapor de ácido acético forma dímeros em boa parte, e a densidade real em relação ao H2 pode passar de 30 em temperaturas baixas. A prova trata o vapor como monomérico e ideal. Nenhuma das duas resoluções comenta isso, mas a questão não pede.%
\item%
\textit{(solution)} A resolução usa M(H2) = 2 e massas inteiras (C = 12, O = 16, H = 1), em vez dos valores da tabela da prova (1,0079; 12,011; 15,999). O gabarito faz o mesmo. A resolução ainda confere com os valores da tabela (M $\approx$ 60,05; razão $\approx$ 29,8 $\approx$ 30), e isso não muda a resposta.%
\end{itemize}%
\section*{Questão 7}%
\href{https://poliedroresolve.sistemapoliedro.com.br/vestibulares/ime/2026/ime-2-fase-dia-3-25-10-2025-2-fase-dia-3/dissertativa-7-quimica-2-fase-dia-3-ime-2026}{Resolução do gabarito}\par\smallskip%
\begin{longtable}{@{}p{0.06\textwidth} p{0.33\textwidth} p{0.33\textwidth} p{0.1\textwidth} p{0.08\textwidth}@{}}\toprule Item & Resolução & Gabarito & Veredito & Certo \\ \midrule
Q7 -- número mínimo de placas solares em 1 h & 42 placas ($\Delta$H\_fus(Cu) = 12,9648 kJ/mol; $\Delta$H global = +49,93 kJ por mol de malaquita, ou seja, por 2 mol de Cu; Q $\approx$ 83,8 MJ; N $\geq$ 41,9, então 42) & 42 placas (T\_fus = 1296,48 K; $\Delta$H\_fus = +12,9648 kJ/mol; $\Delta$H global = +49,9296 kJ por 2 mol de Cu; Q = 83,82 MJ; 83,82/2 = 41,91, então 42) & igual & ambos \\
\bottomrule\end{longtable}%
\begin{itemize}%
\item%
\textbf{Q7 -- número mínimo de placas solares em 1 h)} Refiz as contas. $\Delta$G\_fus = G\_L - G\_S = 12964,8 - 10,0$\cdot$T, o que dá $\Delta$H\_fus = 12964,8 J/mol e $\Delta$S\_fus = 10 J/(mol$\cdot$K), em acordo com a folha de dados. Reação global: Cu2CO3(OH)2 + C $\to$ 2Cu(l) + H2O(g) + 2CO2(g). $\Delta$H = 2(12,9648) - 242 - 2(394) + 1054 = +49,9296 kJ. n\_Cu = 213360/63,546 = 3357,57 mol. Q = (49,9296/2)$\cdot$3357,57 = 83 821 kJ = 83,82 MJ. N = 83,82/2 = 41,91, então 42 placas. Com M = 63,5 g/mol, N = 41,94, que também dá 42. O gabarito tira $\Delta$H\_fus impondo $\Delta$G = 0 em T\_f = 1296,48 K e usando $\Delta$S = 10 da folha. A resolução identifica o termo constante e o coeficiente de T. Os dois caminhos são equivalentes e chegam ao mesmo número.%
\end{itemize}%
\paragraph{Pontos de dúvida}%
\begin{itemize}%
\item%
\textit{(statement)} O enunciado diz que as expressões de G estão em kJ/mol. Nessa unidade, $\Delta$H\_fus seria 12 964,8 kJ/mol e $\Delta$S\_fus seria 10 kJ/(mol$\cdot$K), o que é absurdo e não bate com $\Delta$S\_fus = 10,0 J/(mol$\cdot$K) da folha de dados. A resolução e o gabarito tratam as expressões como J/mol, e o gabarito faz isso sem dizer (1296,48$\cdot$10 = 12 964,8, escrito como 12,9648 kJ). Quem lesse ao pé da letra em kJ/mol chegaria a um número de placas muito diferente. Vale ver se a banca aceitou essa leitura ou anulou o item.%
\item%
\textit{(statement)} A temperatura de fusão que sai das expressões é T\_f = 1296,48 K. Ela não coincide com o limite superior da faixa dada (1358 K, que é o ponto de fusão real do Cu). Isso não afeta o cálculo, porque $\Delta$H é tomado como constante (item iii), mas é uma inconsistência interna dos dados.%
\item%
\textit{(solution)} Há hipóteses implícitas, comuns à resolução e ao gabarito. A energia que as placas precisam fornecer é $\Delta$H global, e não $\Delta$G. O aproveitamento é de 100\% (sem perdas na conversão elétrica para térmica). Também não se conta o calor sensível para aquecer reagentes e produtos até a temperatura de fusão. A resolução declara a hipótese de aproveitamento total. Nenhum dos dois discute a alternativa $\Delta$G, que o enunciado não sugere.%
\item%
\textit{(key)} Na redação, o gabarito chama as placas de 'células voltaicas' e escreve '$\Delta$H\_f' para a entalpia de fusão, o que se confunde com entalpia de formação. São só imprecisões de notação, sem efeito no resultado.%
\end{itemize}%
\section*{Questão 8}%
\href{https://poliedroresolve.sistemapoliedro.com.br/vestibulares/ime/2026/ime-2-fase-dia-3-25-10-2025-2-fase-dia-3/dissertativa-8-quimica-2-fase-dia-3-ime-2026}{Resolução do gabarito}\par\smallskip%
\begin{longtable}{@{}p{0.06\textwidth} p{0.33\textwidth} p{0.33\textwidth} p{0.1\textwidth} p{0.08\textwidth}@{}}\toprule Item & Resolução & Gabarito & Veredito & Certo \\ \midrule
Questão 8 (deficiência percentual de elétrons de valência do {[}B12H12{]}2-) & $\approx$ 40,5\% (50 elétrons reais contra 84 esperados: 42 ligações 2c--2e, sendo 30 B--B e 12 B--H; déficit de 34 elétrons, 34/84) & $\approx$ 40,5\% ((84 - 50)/84 $\times$ 100\%; 30 arestas B--B pela relação de Euler + 12 B--H = 42 ligações $\times$ 2 = 84 elétrons) & igual & ambos \\
\bottomrule\end{longtable}%
\begin{itemize}%
\item%
\textbf{Questão 8 (deficiência percentual de elétrons de valência do {[}B12H12{]}2-))} Valor, método e definição batem. Conferi as contas: elétrons reais 12$\cdot$3 + 12$\cdot$1 + 2 = 50. O icosaedro tem 30 arestas (Euler: 12 + 20 = A + 2). Esperados: (30 + 12)$\cdot$2 = 84. D = 34/84 = 40,48\% $\approx$ 40,5\%. Reli o enunciado no PDF original (pág. 4) e ele confere com a transcrição. A figura mostra as 30 arestas, incluindo as tracejadas, e um H terminal em cada vértice. A resolução acrescenta a verificação 3F = 2A, a analogia com o B2H6 (25\%) e a regra de Wade (13 pares de esqueleto = n+1). Nada disso contradiz o gabarito.%
\end{itemize}%
\paragraph{Pontos de dúvida}%
\begin{itemize}%
\item%
\textit{(statement)} O enunciado não define a fórmula da 'deficiência percentual'. Resolução e gabarito dividem pelo número esperado (34/84 $\approx$ 40,5\%), que é a leitura mais natural de 'em relação ao esperado'. Há outras leituras possíveis. Dividindo pelo número real, o resultado seria 34/50 = 68\%. Contando só o esqueleto B12 (30 ligações B--B, 60 elétrons esperados), seria 34/60 $\approx$ 56,7\%. Uma correção que aceite essas variantes, desde que bem justificadas, mudaria a nota de quem as usou, mas não a desta resolução.%
\item%
\textit{(statement)} A hipótese de que 'só ligações 2c--2e' quer dizer uma ligação em cada aresta do icosaedro deixaria cada B com 6 ligações (5 B--B + 1 B--H), o que não é típico para o boro. Resolução e gabarito adotam essa leitura geométrica, e ela é a que a figura sugere. Fica registrado como hipótese implícita, sem impacto no resultado esperado.%
\end{itemize}%
\section*{Questão 9}%
\href{https://poliedroresolve.sistemapoliedro.com.br/vestibulares/ime/2026/ime-2-fase-dia-3-25-10-2025-2-fase-dia-3/dissertativa-9-quimica-2-fase-dia-3-ime-2026}{Resolução do gabarito}\par\smallskip%
\begin{longtable}{@{}p{0.06\textwidth} p{0.33\textwidth} p{0.33\textwidth} p{0.1\textwidth} p{0.08\textwidth}@{}}\toprule Item & Resolução & Gabarito & Veredito & Certo \\ \midrule
a & {[}CO2{]} $\approx$ 1,8$\times$10\^{}{-}5 mol/L (p\_CO2 = 5,45$\times$10\^{}{-}4 atm; lei de Henry com k\_H = 3,3$\times$10\^{}{-}2 mol/(atm$\cdot$L)) & {[}CO2{]} $\approx$ 1,8$\cdot$10\^{}{-}5 mol/L (P\_CO2 = 545$\cdot$10\^{}{-}6 atm; S = k\_H$\cdot$P) & igual & ambos \\
b & $\Delta$G$^\circ$ = +36,3 kJ/mol; K $\approx$ 4,3$\times$10\^{}{-}7 (ln K $\approx$ -14,68); {[}H+{]} $\approx$ 2,8$\times$10\^{}{-}6 mol/L; pH $\approx$ 5,56 $\approx$ 5,6 & $\Delta$G$^\circ$ = +36,3 kJ/mol; K = e\^{}(-14,68); 2pH $\approx$ 11,11; pH $\approx$ 5,6 & igual & ambos \\
c & O pH diminui: $\Delta$H$^\circ$ = +7,6 kJ/mol (ionização endotérmica), K cresce com T (van't Hoff) e, com {[}CO2{]} fixa, {[}H+{]} = $\surd$(K{[}CO2{]}) aumenta & O pH diminui: $\Delta$H$^\circ$ = +7,6 kJ/mol, reação endotérmica; com o aumento de T, K e {[}H+{]} aumentam e o pH cai & igual & ambos \\
\bottomrule\end{longtable}%
\begin{itemize}%
\item%
\textbf{a)} Recalculei: 3,3$\times$10\^{}{-}2 $\times$ 5,45$\times$10\^{}{-}4 = 1,7985$\times$10\^{}{-}5 $\approx$ 1,8$\times$10\^{}{-}5 mol/L. Os dois tomam P\_atm = 1 atm e usam a constante de Henry dada para 300 K como se valesse a 25 $^\circ$C. O enunciado conferido no PDF (p. 4) bate com a transcrição.%
\item%
\textbf{b)} $\Delta$G$^\circ$ = -586,8 + 237,1 + 386,0 = +36,3 kJ/mol. ln K = -36300/(8,3$\cdot$298) = -14,676, então log K $\approx$ -6,37 e K $\approx$ 4,2--4,3$\times$10\^{}{-}7. Com {[}H+{]}$^2$ = K$\cdot${[}CO2{]}, o pH dá 5,559; com 1,8$\times$10\^{}{-}5 arredondado, 2pH = 11,115 (o gabarito escreve 11,11). Se incluo Kw ({[}H+{]}$^2$ = K$\cdot$c + Kw), o pH fica em 5,559, sem mudança. A resolução se justifica um pouco melhor: verifica que {[}OH-{]} e a 2ª ionização são desprezíveis.%
\item%
\textbf{c)} $\Delta$H$^\circ$ = -692,0 + 413,8 + 285,8 = +7,6 kJ/mol > 0. Pela equação de van't Hoff, K aumenta com T; como {[}CO2{]} fica constante, {[}H+{]} aumenta e o pH diminui. A estimativa de ordem de grandeza da resolução está certa: de 298 K para 308 K, ln(K2/K1) $\approx$ 0,10, o que dá $\Delta$pH $\approx$ -0,02. A observação sobre a solubilidade do gás ($\Delta$H\_dissolução = -19,8 kJ/mol) está correta e corretamente posta de lado pelo enunciado.%
\end{itemize}%
\paragraph{Pontos de dúvida}%
\begin{itemize}%
\item%
\textit{(statement)} A constante de Henry vem da folha de dados para 300 K, mas o item (a) pede 25 $^\circ$C (298 K). A resolução e o gabarito a usam sem correção; a resolução declara essa hipótese. O enunciado também não dá a pressão atmosférica total, e os dois assumem P = 1 atm.%
\item%
\textit{(key)} Erro de digitação no gabarito, item (c): escreve $\Delta$H\_f$^\circ${[}CO3\^{}{-}(aq){]} em vez de $\Delta$H\_f$^\circ${[}HCO3\^{}{-}(aq){]}. O valor usado (-692,0) é o do HCO3\^{}{-}, então o resultado não muda. No item (b), a equação da reação também saiu truncada ("H\^{}+ + HCO").%
\item%
\textit{(key)} O gabarito escreve 2pH = 11,11, mas a soma -0,954 - 0,301 + 6 + 6,37 dá 11,115. É só arredondamento: pH $\approx$ 5,56 $\approx$ 5,6 nos dois.%
\item%
\textit{(solution)} A resolução dá K $\approx$ 4,3$\times$10\^{}{-}7, obtido de 10\^{}{-}6,37. O valor direto e\^{}(-14,676) = 4,23$\times$10\^{}{-}7 $\approx$ 4,2$\times$10\^{}{-}7. A diferença é só de arredondamento e não afeta o pH (5,56).%
\item%
\textit{(statement)} No item (c), os dois desprezam a variação de $\Delta$H$^\circ$ com T e o efeito de Kw, o que é razoável para um raciocínio qualitativo. A resolução justifica por que Kw não importa.%
\end{itemize}%
\section*{Questão 10}%
\href{https://poliedroresolve.sistemapoliedro.com.br/vestibulares/ime/2026/ime-2-fase-dia-3-25-10-2025-2-fase-dia-3/dissertativa-10-quimica-2-fase-dia-3-ime-2026}{Resolução do gabarito}\par\smallskip%
\begin{longtable}{@{}p{0.06\textwidth} p{0.33\textwidth} p{0.33\textwidth} p{0.1\textwidth} p{0.08\textwidth}@{}}\toprule Item & Resolução & Gabarito & Veredito & Certo \\ \midrule
a & $^2$$^2$$^{7}$$_{8}$$_{9}$X $\to$ $^2$$^2$$^{7}$$_{9}$$_{0}$J + $\beta$$^{-}$ (+ antineutrino); $^2$$^2$$^{7}$$_{9}$$_{0}$J $\to$ $^2$$^2$$^3$$_{8}$$_{8}$L + $\alpha$; $^2$$^2$$^3$$_{8}$$_{8}$L $\to$ $^2$$^{1}$$^{9}$$_{8}$$_{6}$Z + $\alpha$; $^2$$^{1}$$^{9}$$_{8}$$_{6}$Z $\to$ $^2$$^{1}$$^{5}$$_{8}$$_{4}$T + $\alpha$. Global: X $\to$ T + 3$\alpha$ + $\beta$$^{-}$ + antineutrino, com o balanço de A e Z conferido em cada etapa. & $^2$$^2$$^{7}$$_{8}$$_{9}$X $\to$ $^2$$^2$$^{7}$$_{9}$$_{0}$J + $^{0}$?$_{1}$$\beta$; $^2$$^2$$^{7}$$_{9}$$_{0}$J $\to$ $^2$$^2$$^3$$_{8}$$_{8}$L + $^{4}$$_{2}$$\alpha$; $^2$$^2$$^3$$_{8}$$_{8}$L $\to$ $^2$$^{1}$$^{9}$$_{8}$$_{6}$Z + $^{4}$$_{2}$$\alpha$; $^2$$^{1}$$^{9}$$_{8}$$_{6}$Z $\to$ $^2$$^{1}$$^{5}$$_{8}$$_{4}$T + $^{4}$$_{2}$$\alpha$ & igual & ambos \\
b & Família do actínio (4n+3): 227, 223, 219 e 215 deixam resto 3 na divisão por 4, e esse resto não muda porque $\alpha$ tira 4 de A e $\beta$$^{-}$ não altera A. & Série do actínio (4n+3): 227, 223, 219 e 215 deixam resto 3 na divisão por 4; X (Z = 89, A = 227) é o actínio. & igual & ambos \\
c & $\lambda$ = ln2/t$\frac12$ $\approx$ 0,173 s$^{-}$$^{1}$ (declarado $\approx$ 0,17 s$^{-}$$^{1}$); $\tau$ = 1/$\lambda$ = t$\frac12$/ln2 $\approx$ 5,77 s (declarado $\approx$ 5,8 s) & $\lambda$ = (0,301/0,434)/4 $\approx$ 0,173 s$^{-}$$^{1}$; $\tau$ = 1/$\lambda$ $\approx$ 5,77 s & equivalente & ambos \\
\bottomrule\end{longtable}%
\begin{itemize}%
\item%
\textbf{a)} As quatro etapas e as partículas emitidas (1 $\beta$$^{-}$ e depois 3 $\alpha$) são as mesmas. A resolução vai um pouco além: inclui o antineutrino, que tem A = 0 e carga 0 e por isso não mexe no balanço, e escreve a equação global somada, que confere: 227 = 215 + 12 + 0 e 89 = 84 + 6 - 1. Assim, ela atende a qualquer uma das leituras de "equação estequiométrica da série", seja por etapa, seja global.%
\item%
\textbf{b)} A conclusão e o critério são os mesmos. A resolução ainda explica por que o resto se mantém ao longo da série, o que responde de fato ao "demonstre". Ela também identifica X como $^2$$^2$$^{7}$Ac.%
\item%
\textbf{c)} Refiz as contas: ln2 $\approx$ 0,301/0,434 = 0,6935, então $\lambda$ = 0,6935/4 = 0,1734 s$^{-}$$^{1}$ e $\tau$ = 4/0,6935 = 5,767 s. Com ln2 exato dá $\lambda$ = 0,1733 s$^{-}$$^{1}$ e $\tau$ = 5,771 s. Os números da resolução, 0,17 s$^{-}$$^{1}$ e 5,8 s, diferem do gabarito só pelo arredondamento final, e a própria resolução também mostra 0,173 s$^{-}$$^{1}$ e 5,77 s. As unidades estão no SI (s$^{-}$$^{1}$ e s), como o enunciado pede.%
\end{itemize}%
\paragraph{Pontos de dúvida}%
\begin{itemize}%
\item%
\textit{(key)} Segundo a nota do gabarito, o enunciado traria massa 277 para X e J, e o Poliedro teria usado 227 sem comentar. No PDF original desta pasta (página 5), tanto na imagem quanto no texto extraído com pdftotext, aparece 227 para X e J. Para esta versão da prova não há erro de digitação, e 227 é o único valor coerente com a série: 227 $\to$ 223 por emissão $\alpha$, e Z = 89 é o actínio. Se existir outra versão da prova com 277, vale confirmar com o professor.%
\item%
\textit{(solution)} Na resposta final, a resolução arredonda para $\lambda$ $\approx$ 0,17 s$^{-}$$^{1}$ e $\tau$ $\approx$ 5,8 s, enquanto o gabarito traz 0,173 s$^{-}$$^{1}$ e 5,77 s. É só arredondamento, e a resolução também mostra os valores com três algarismos no desenvolvimento.%
\item%
\textit{(statement)} "Equação estequiométrica da série" pode ser entendida como as equações de cada etapa (leitura do gabarito) ou como uma equação global. A resolução apresenta as duas formas, então não perde nota em nenhuma leitura. A atmosfera ambiente e as CNTP são distratores, porque $\lambda$ não depende das condições externas, como a resolução observa.%
\end{itemize}%
\end{document}